\section{Related Work}

Detailed related work is discussed in Appendix~\ref{app:related-work}.


\paragraph{LLMs for Data Science:} Specialized benchmarks evaluate LLMs in data science. DS-1000 \citep{10.5555/3618408.3619164}, ARCADE \citep{yin-etal-2023-natural}, DataSciBench \citep{zhang2025datascibenchllmagentbenchmark}, and DSEval \citep{zhang-etal-2024-benchmarking-data} focus on translating natural language instructions into correct implementations, unlike broader programming benchmarks such as SWE-Bench \citep{jimenez2024swebench}, ML-Bench \citep{tang2025mlbench}, and BigCodeBench \citep{zhuo2025bigcodebench}. Most assume relevant data files are pre-specified, though DSBench \citep{jing2025dsbench} and BLADE \citep{gu-etal-2024-blade} address multi-step reasoning, while ScienceAgentBench \citep{chen2025scienceagentbench} and BixBench \citep{mitchener2025bixbenchcomprehensivebenchmarkllmbased} focus on domain knowledge integration. These benchmarks largely overlook the challenge of discovering relevant data in large, heterogeneous repositories—a gap addressed by KramaBench \citep{lai2025kramabenchbenchmarkaisystems}, which explicitly evaluates data discovery. Building on this, we study how agents can autonomously identify and leverage correct data sources for end-to-end analysis. Applications of LLMs in data science have evolved from single-turn code generation to interactive, tool-augmented agents using models specialized for code, including GPT \citep{10.5555/3495724.3495883}, CodeGen \citep{rubavicius2025conversationalcodegenerationcase}, StarCoder \citep{li2023starcoder}, and Code Llama \citep{codellama}. While few-shot prompting \citep{10.5555/3495724.3495883} remains effective, state-of-the-art methods adopt agentic or multi-agentic frameworks combining reasoning with tool use. ReAct \citep{yao2023react} interleaves reasoning with actions, extended to execution environments \citep{chen2018executionguided}. Toolformer \citep{schick2023toolformer} and Gorilla \citep{patil2024gorilla} train LLMs to call APIs for tasks requiring specialized libraries. Self-correction frameworks like Self-Debug \citep{chen2024teaching} and Reflexion \citep{shinn2023reflexion} refine code using execution feedback. To improve reliability, many systems integrate RAG \citep{10.5555/3495724.3496517, salemi2025planandrefinediversecomprehensiveretrievalaugmented, 10.1145/3731120.3744584, 10.1145/3626772.3657733} to access documentation or code examples, reducing hallucinations and ensuring up-to-date library use. Multi-agent master–slave frameworks, such as AutoKaggle \citep{li2025autokaggle}, further demonstrate promising results in this domain.


\paragraph{Blackboard Systems:}

The blackboard system is a architectural model from classical AI, developed for problems that require incremental and opportunistic reasoning. It was implemented in the Hearsay-II speech understanding \citep{10.1145/356810.356816} and is characterized by three components: (1) a global, hierarchical data structure (the blackboard) that maintains the current state of the solution; (2) independent specialist modules, known as knowledge sources, which monitor the blackboard and contribute partial solutions; and (3) a control mechanism that opportunistically determines which knowledge source to activate next \citep{Nii_1986}. Following successful applications in domains such as sonar interpretation with the HASP/SIAP system \citep{Nii_Feigenbaum_Anton_1982}, the architecture evolved to incorporate more sophisticated control strategies. Inspired by this, we adapt the blackboard architecture for LLM-based multi-agent communication. Instead of a central controller assigning tasks, agents operate autonomously and respond to requests posted on the blackboard when they can contribute without being forced to participate. The main agent then integrates the information provided by the sub-agents to solve the problem.